\chapter{Flows and Matchings}\label{ch:flows}

\begin{center}\itshape
A matching is maximum exactly when no alternating path joins two free vertices;\\
the whole difficulty is to find such a path when odd cycles hide it.\\
\upshape(Berge's theorem and Edmonds' problem, in our words)
\end{center}

In 1965 the \emph{Canadian Journal of Mathematics} published a paper by Jack Edmonds, a young mathematician at the National Bureau of Standards in Washington,
with an unusual title: \emph{Paths, trees, and flowers}. It solved a concrete problem, finding a largest matching in an arbitrary graph, and it did so with
an algorithm whose key move was to shrink odd cycles, which Edmonds called \emph{blossoms}, into single vertices. But it also contained a short digression that
changed computer science. There is an obvious finite algorithm for the problem, Edmonds observed, which tries everything, and its work grows exponentially with the
size of the graph; by contrast, he called an algorithm \emph{good} if its work is bounded by a polynomial in the size of the input, and he claimed, as a mathematical
result, that a good algorithm for maximum matching exists. (This is our paraphrase of his Section~2.) Polynomial time as the meaning of ``efficient'' starts there, and Tome~C3
will build on it.

Tome~C1, Chapter~6, proved the theorems: Berge's augmenting path criterion, Hall and K\H onig, the max-flow min-cut theorem with the integrality theorem, Tutte's
perfect matching theorem. This chapter makes them fast. We bound the number of augmentations when the shortest augmenting path is always chosen (Edmonds and Karp), group
augmentations into phases (Dinic, and Hopcroft and Karp for bipartite matching), prove Edmonds' blossom algorithm completely, and then add costs: minimum-cost flows by successive
shortest paths, whose special case is the Hungarian algorithm for the assignment problem.

\begin{think}
In a bipartite graph, a breadth-first search from a free vertex, alternating between non-matching and matching edges, finds an augmenting path whenever one exists.
Try the same search on a triangle $xyz$ with the matching $\{yz\}$ and an extra vertex $w$ joined to $z$ only, starting from $x$. What goes wrong?
\end{think}

\section{Residual networks and shortest augmenting paths}\label{sec:edmondskarp}

We keep the notation of Tome~C1: a network $D=(V,A)$ with source $s$, sink $t$, capacities $c\ge0$, flows $f$ of value $|f|$; $n=|V|$ and $m=|A|$. We assume that no two arcs are
opposite (subdividing one arc of each opposite pair changes nothing), so that the following is well defined. The \emph{residual network} $D_f$ has, for every arc $uv$, the arc $uv$ with
\emph{residual capacity} $r(uv)=c(uv)-f(uv)$ if this is positive, and the reverse arc $vu$ with $r(vu)=f(uv)$ if this is positive. An augmenting path of Tome~C1 is exactly a path from $s$ to $t$ in $D_f$; augmenting by
its bottleneck $\delta=\min r$ adds $\delta$ to forward arcs and subtracts it from reversed ones, and at least one arc of the path, a \emph{critical} arc, disappears from the residual network.

Tome~C1, Theorem~6.6, showed that augmenting along arbitrary paths may never stop when capacities are irrational, and with integer capacities may take as many steps as the value of the flow
(Exercise~\ref{ex:ffbad}). In 1972 Edmonds and Karp proposed to always augment along a \emph{shortest} path, found by breadth-first search (Theorem~\ref{thm:bfs}). Write $d_f(v)$ for the distance from $s$ to $v$ in $D_f$
(in number of arcs).

\begin{lemma}\label{lem:monotonedist}
If $f'$ is obtained from $f$ by augmenting along a shortest path, then $d_{f'}(v)\ge d_f(v)$ for every vertex $v$.
\end{lemma}
\begin{proof}
Suppose not, and among the vertices with $d_{f'}(v)<d_f(v)$ take one with $d_{f'}(v)$ least; $v\ne s$. Let $u$ precede $v$ on a shortest path from $s$ in $D_{f'}$, so $d_{f'}(u)=d_{f'}(v)-1$, and by minimality $d_f(u)\le d_{f'}(u)$.
If the arc $uv$ is in $D_f$, then $d_f(v)\le d_f(u)+1\le d_{f'}(u)+1=d_{f'}(v)$, a contradiction. Otherwise $uv$ appeared during the augmentation, which is only possible if the reverse arc $vu$ was on the shortest path used, so
$d_f(u)=d_f(v)+1$. Then $d_f(v)=d_f(u)-1\le d_{f'}(u)-1=d_{f'}(v)-2$, a contradiction again.
\end{proof}

\begin{theorem}[Edmonds and Karp, 1972; Dinic, 1970]\label{thm:edmondskarp}
Augmenting along shortest paths finds a maximum flow after at most $nm$ augmentations, whatever the capacities, in time $O(nm^2)$.
\end{theorem}
\begin{proof}
Each residual arc is an arc of $D$ or its reverse: there are at most $2m$ of them. Suppose $uv$ is critical in an augmentation from $f$; it was on a shortest path, so $d_f(v)=d_f(u)+1$, and it disappears.
It can reappear only after an augmentation, from some later flow $g$, that uses $vu$, so that $d_g(u)=d_g(v)+1\ge d_f(v)+1=d_f(u)+2$ by Lemma~\ref{lem:monotonedist}. So between two occasions where $uv$ is critical, the distance of $u$ grows by at least $2$.
Since $u$ is reachable at those moments, $d(u)\le n-2$ (as $u\ne t$), so each residual arc is critical at most $n/2$ times. Every augmentation has a critical arc, hence there are at most $2m\cdot n/2=nm$ augmentations, and each costs $O(m)$
(a breadth-first search, then the update). When no path remains, the flow is maximum by Tome~C1, Theorem~6.6.
\end{proof}

\section{Phases: Dinic, and Hopcroft and Karp}\label{sec:dinic}

Yefim Dinitz, a student in Moscow, found a closely related algorithm in 1970 (his name was transliterated ``Dinic'' in the translated paper, and the algorithm kept that name). Its idea is to do all augmentations
of the same length at once. The \emph{level graph} $L_f$ keeps only the arcs $uv$ of $D_f$ with $d_f(v)=d_f(u)+1$; every shortest path of $D_f$ from $s$ to $t$ lies in it, and it is acyclic. A \emph{blocking flow} in $L_f$ is a flow
(with the residual capacities as capacities) such that every path from $s$ to $t$ in $L_f$ contains a saturated arc. \emph{Dinic's algorithm} repeats \emph{phases}: compute $L_f$ by breadth-first search, find a blocking flow $g$ in it,
and add $g$ to $f$ (adding on arcs of $D$, subtracting on reversed ones).

\begin{theorem}[Dinic, 1970]\label{thm:dinic}
Each phase increases $d_f(t)$, so there are at most $n-1$ phases. A blocking flow can be found in time $O(nm)$, hence a maximum flow in time $O(n^2m)$.
\end{theorem}
\begin{proof}
Let $d=d_f$ and let $f'$ be the flow after the phase. An arc of $D_{f'}$ is either an arc of $D_f$, which satisfies $d(v)\le d(u)+1$, or the reverse of an arc of $L_f$, which satisfies $d(v)=d(u)-1$. Along a path from $s$ to $t$ in $D_{f'}$,
$d$ increases by at most $1$ per arc, so the path has at least $d(t)$ arcs, with equality only if every arc satisfies $d(v)=d(u)+1$ and is therefore an arc of $L_f$ that kept a positive residual capacity, that is, an arc of $L_f$ not saturated by $g$.
Such a path would contradict the blocking property. So $d_{f'}(t)>d_f(t)$; since $d_f(t)\le n-1$ while $t$ is reachable, there are at most $n-1$ phases.

To find a blocking flow, keep for each vertex a pointer to the next unexamined arc of $L_f$ leaving it, and walk from $s$ following pointers (\emph{advance}). If the walk reaches $t$, augment along it by its bottleneck,
delete the saturated arcs, and restart from $s$. If it gets stuck at a vertex $v\ne t$ whose arcs are all examined, delete the arc that led to $v$ and step back (\emph{retreat}); when $s$ itself is stuck, stop. A deleted arc is either saturated
or leads to a vertex from which, by induction on the order of deletions, every remaining path to $t$ is blocked; so at the end every path from $s$ to $t$ in $L_f$ contains a saturated arc. Each augmentation deletes an arc and costs $O(n)$, the length
of the path; each retreat deletes an arc and cancels one advance; so there are at most $m$ augmentations and $m$ retreats, and at most $m+nm$ advances. The phase costs $O(nm)$.
\end{proof}

On \emph{unit} networks, where all capacities are $1$, a phase costs only $O(m)$ and the number of phases is much smaller (Exercise~\ref{ex:dinicunit}). The most important unit network comes from matchings: for a bipartite graph with sides $X$ and $Y$,
add a source joined to $X$ and a sink joined from $Y$, and orient the edges from $X$ to $Y$; integral flows are matchings (Tome~C1, Chapter~6). Hopcroft and Karp found the algorithm directly in 1973, in the language of matchings.

Let $M$ be a matching of a graph $G$ (bipartite or not). Recall from Tome~C1, Theorem~6.1 (Berge), that an $M$-augmenting path joins two free vertices and alternates between edges outside and inside $M$, and that $M\triangle P$ is then a matching
with one more edge. The \emph{length} of a path is its number of edges.

\begin{lemma}\label{lem:symdiff}
If $M$ and $N$ are matchings with $|N|>|M|$, then $M\triangle N$ contains at least $|N|-|M|$ vertex-disjoint $M$-augmenting paths.
\end{lemma}
\begin{proof}
Every vertex meets at most one edge of each matching, so $M\triangle N$ has maximum degree $2$: its components are paths and cycles whose edges alternate between $M$ and $N$. Cycles are even and contain as many edges of $M$ as of $N$;
a path has at most one more edge of one kind. Since $N$ has $|N|-|M|$ more edges than $M$ in $M\triangle N$, at least $|N|-|M|$ components are paths with one more $N$-edge; their ends meet no $M$-edge of the path, and none outside it either,
since an $M$-edge at an end would belong to $M\triangle N$ or to $M\cap N$, impossible as the end meets an edge of $N$ in the path. So they are $M$-augmenting.
\end{proof}

\begin{lemma}\label{lem:hkphase}
Let $\ell$ be the length of a shortest $M$-augmenting path, let $P_1,\dots,P_k$ be a \emph{maximal} family of vertex-disjoint $M$-augmenting paths of length $\ell$, and let $M'=M\triangle P_1\triangle\dots\triangle P_k$. Then every $M'$-augmenting path has length
at least $\ell+2$.
\end{lemma}
\begin{proof}
Let $P$ be $M'$-augmenting, and $M''=M'\triangle P$, of size $|M|+k+1$. By Lemma~\ref{lem:symdiff}, $M\triangle M''$ contains $k+1$ vertex-disjoint $M$-augmenting paths, each of length at least $\ell$, so $|M\triangle M''|\ge(k+1)\ell$. But
$M\triangle M''=(P_1\cup\dots\cup P_k)\triangle P$, so with $Q=P_1\cup\dots\cup P_k$,
\[
(k+1)\ell\le|Q|+|P|-2|Q\cap P|=k\ell+|P|-2|Q\cap P|,
\]
that is, $|P|\ge\ell+2|Q\cap P|$. If $P$ shares no vertex with $Q$, its edges and ends have the same status for $M$ as for $M'$, so it is $M$-augmenting and vertex-disjoint from the $P_i$; by maximality of the family, $|P|\ne\ell$, so $|P|>\ell$, and in fact
$|P|\ge\ell+2$ because augmenting paths, which start and end with non-matching edges, have odd length. If $P$ shares a vertex $v$ with some $P_i$, then $v$ is matched in $M'$ by an edge of $P_i$, so $v$ is not an end of $P$,
and $P$, alternating, contains the $M'$-edge at $v$: $|Q\cap P|\ge1$ and $|P|\ge\ell+2$.
\end{proof}

\begin{theorem}[Hopcroft and Karp, 1973]\label{thm:hopcroftkarp}
A maximum matching of a bipartite graph is found in time $O(m\sqrt n)$, by phases each of which augments along a maximal family of vertex-disjoint shortest augmenting paths.
\end{theorem}
\begin{proof}
By Lemma~\ref{lem:hkphase}, the shortest augmenting length increases by at least $2$ per phase, so after $\lceil\sqrt n\rceil$ phases it exceeds $\sqrt n$. Let $M$ be the matching at that moment and $M^*$ a maximum matching. By Lemma~\ref{lem:symdiff},
$M\triangle M^*$ contains $|M^*|-|M|$ vertex-disjoint $M$-augmenting paths, each with more than $\sqrt n$ vertices, so $|M^*|-|M|<\sqrt n$. Every later phase adds at least one edge, so at most $\sqrt n$ more phases follow: $O(\sqrt n)$ phases in all.

A phase costs $O(m)$. A breadth-first search from all free vertices of $X$ at once, using non-matching edges from $X$ to $Y$ and matching edges from $Y$ to $X$, labels every vertex with its alternating distance and stops at the first level containing a free vertex of $Y$;
this gives $\ell$ and the layered graph of all shortest augmenting paths. Then a depth-first search from each free vertex of $X$, in the layered graph, looks for a path to a free vertex of $Y$ at level $\ell$, deleting every vertex it visits, whether it succeeds or not.
Deleted vertices that led nowhere lead nowhere later, and those on a path found are used; so the paths found are vertex-disjoint and the family is maximal, and each edge is examined a constant number of times.
\end{proof}

\section{Blossoms}\label{sec:blossom}

In a general graph the search of the bipartite case fails, as the opening question shows: from the free vertex $x$ of the triangle, the alternating search reaches $y$ and $z$ by non-matching edges, and then the matching edge $yz$ leads from each
to the other, already visited; the augmenting path $x,y,z,w$ is missed because $z$ is reached at the wrong \emph{parity}. Odd cycles are the culprit. Edmonds' algorithm builds alternating trees and, when it meets an odd cycle, shrinks it.

Fix a matching $M$ of $G$. An \emph{alternating forest} $F$ has the free vertices as roots, and alternates as one goes down: each vertex at even depth (\emph{even}, or outer) other than a root is joined to its parent by an $M$-edge, each vertex at odd depth
(\emph{odd}, or inner) is joined to its parent by a non-$M$ edge and has exactly one child, its mate. So the path from a root to any vertex of $F$ is alternating, and starts with a non-$M$ edge. Initially $F$ consists of the free vertices. Repeatedly
look at an edge $vw$ with $v$ even:
\begin{enumerate}
\item if $w\notin F$, then $w$ is matched (free vertices are in $F$) and its mate $w'$ is not in $F$ either; \emph{grow}: add $w$ as an odd child of $v$, and $w'$ as an even child of $w$;
\item if $w$ is even and in another tree, the path from the root of $v$ to $v$, the edge $vw$, and the path from $w$ to its root form an $M$-augmenting path: \emph{augment};
\item if $w$ is even and in the same tree, the tree paths from $v$ and $w$ to their lowest common ancestor $b$, with the edge $vw$, form a cycle $B$ of odd length $2k+1$: a \emph{blossom}, with \emph{base} $b$, which is even;
\item if $w$ is odd, do nothing.
\end{enumerate}
In case~3, the $2k$ vertices of $B$ other than $b$ are matched by $k$ edges of the cycle, and the tree path from the root $r$ to $b$, the \emph{stem}, is alternating of even length, ending with the $M$-edge at $b$ (if $b\ne r$). Replacing $M$ by $M\triangle\text{stem}$
keeps its size, makes $b$ free and $r$ matched, and leaves the matching inside $B$ unchanged. Then we \emph{shrink}: $G/B$ has the vertices outside $B$ and one new vertex, also called $B$; each edge with one end in $B$ gets that end replaced by the new vertex, and
edges inside $B$ disappear. Since $b$ is free and no other vertex of $B$ is matched outside $B$, the matching $M/B$ (the edges of $M$ not inside $B$) leaves $B$ free and has $k$ edges fewer than $M$. Figure~\ref{fig:blossom} shows the operation.

\begin{figure}[ht]
\centering
\fig{ch06-blossom}
\caption{Left: an alternating tree rooted at $r$; matched edges are thick, even vertices shaded. The edge $df$ (dashed) joins two even vertices of the same tree and closes the blossom $B=b\,c\,d\,f\,e$ with base $b$; the stem is $r\,a\,b$.
Right: after exchanging matched and unmatched edges on the stem, $b$ is free, and $B$ is shrunk to a single free vertex.}
\label{fig:blossom}
\end{figure}

\begin{lemma}[Edmonds, 1965]\label{lem:shrink}
Let $B$ be an odd cycle of length $2k+1$ in $G$ and $M$ a matching such that one vertex $b$ of $B$ is free and the other $2k$ are matched by edges of $B$. Then $M$ is a maximum matching of $G$ if and only if $M/B$ is a maximum matching of $G/B$.
Moreover, a matching $N'$ of $G/B$ can be turned in linear time into a matching of $G$ with $|N'|+k$ edges.
\end{lemma}
\begin{proof}
\emph{Expanding.} In $N'$ the vertex $B$ is matched by at most one edge, which in $G$ ends at some $y\in B$; if $B$ is free in $N'$, let $y=b$. Deleting $y$ from the cycle $B$ leaves a path with $2k$ vertices, which has a perfect matching of $k$
edges. Adding them to $N'$ gives a matching of $G$ of size $|N'|+k$. In particular, if $M/B$ is not maximum, a larger $N'$ gives a matching larger than $|M/B|+k=|M|$, and $M$ is not maximum.

\emph{Conversely}, suppose $M$ is not maximum; by Berge's theorem there is an $M$-augmenting path $P$ in $G$. If $P$ avoids $B$, it is $M/B$-augmenting in $G/B$. Otherwise, $P$ has two distinct free ends, and $b$ is the only free vertex of $B$, so some end $x$ of $P$ is
not in $B$. Follow $P$ from $x$ up to its first vertex $y$ in $B$. This piece starts with a non-$M$ edge and alternates; its last edge enters $B$ at $y$, and is not in $M$, because $y$ is either $b$, which is free, or matched by an edge inside $B$. In $G/B$
the piece becomes a path from the free vertex $x$ to the free vertex $B$ that alternates and ends with a non-matching edge: it is $M/B$-augmenting. So $M/B$ is not maximum.
\end{proof}

\begin{theorem}[Edmonds, 1965]\label{thm:blossom}
The following procedure $\mathrm{Improve}(G,M)$ returns a matching of $G$ with $|M|+1$ edges, or certifies that $M$ is maximum. Grow an alternating forest as above. On case~2, return $M\triangle P$. On case~3, flip the stem, giving $M_1$, shrink the blossom $B$, call
$\mathrm{Improve}(G/B,M_1/B)$, and if it returns a larger matching, expand it to $G$ by Lemma~\ref{lem:shrink} and return it; if it certifies that $M_1/B$ is maximum, certify that $M$ is. If no edge from an even vertex leads to a vertex that is not odd,
certify that $M$ is maximum. Calling $\mathrm{Improve}$ at most $n/2$ times from the empty matching gives a maximum matching of an arbitrary graph, in time $O(n^2m)$.
\end{theorem}
\begin{proof}
We argue by induction on the number of vertices. Case~2 is correct by Berge. In case~3, the flipped matching $M_1$ has the size of $M$, and satisfies the hypothesis of Lemma~\ref{lem:shrink}; $G/B$ has fewer vertices. If the recursive call
returns a matching with $|M_1/B|+1$ edges, its expansion has $|M_1/B|+1+k=|M_1|+1=|M|+1$ edges. If it certifies $M_1/B$ maximum, then $M_1$ is maximum by the lemma, and so is $M$, of the same size.

It remains to treat the case where the forest stops growing: every neighbour of every even vertex is odd. Let $A$ be the set of odd vertices and $E_v$ the set of even ones. Each odd vertex has exactly one child, so $|E_v|=|A|+\varphi$, where $\varphi$ is the number of roots, which
is the number of $M$-free vertices, so $|M|=(n-\varphi)/2$. Now let $N$ be any matching. An even vertex matched by $N$ is matched to one of its neighbours, all in $A$, and each vertex of $A$ is matched to at most one of them; so at least $|E_v|-|A|=\varphi$ vertices are
free in $N$, and $|N|\le(n-\varphi)/2=|M|$. So $M$ is maximum.

For the time: one call scans each edge a constant number of times while growing the forest ($O(n+m)$, with an array recording each vertex's tree, parity and parent), then either augments, flips and shrinks in $O(n+m)$, or stops. Shrinkings nest to depth at most $n/2$,
since each removes at least two vertices, so one improvement costs $O(n(n+m))$, and there are at most $n/2$ improvements: $O(n^2m)$ when $G$ has no isolated vertices.
\end{proof}

The last part of the proof is a little treasure: the stuck forest exhibits a set $A$ whose removal leaves many isolated vertices, which is exactly the kind of certificate in Tutte's theorem (Tome~C1, Theorem~6.10), and it shows that the
Tutte--Berge bound of Exercise~\ref{ex:tutteberge} is attained (when the forest gets stuck without any shrinking; with shrinking the certificate must be expanded, see the sources). Micali and Vazirani (1980) brought the time down to $O(m\sqrt n)$, as for bipartite graphs.

\section{Costs: minimum-cost flows and the assignment problem}\label{sec:mincost}

Now each arc also has a \emph{cost} $\kappa(uv)$ per unit of flow, and the cost of a flow is $\kappa(f)=\sum_a\kappa(a)f(a)$. In the residual network, a forward arc $uv$ costs $\kappa(uv)$ and a reversed one $vu$ costs $-\kappa(uv)$: sending flow back saves its cost.

\begin{lemma}[negative cycles]\label{lem:negcycle}
A flow $f$ has minimum cost among the flows of value $|f|$ if and only if $D_f$ has no cycle of negative cost.
\end{lemma}
\begin{proof}
If $D_f$ has a negative cycle, pushing its bottleneck around it keeps the value and lowers the cost. Conversely, let $f^*$ be any flow with $|f^*|=|f|$. Define on the arcs of $D_f$ the function $h$: on a forward arc $uv$, $h(uv)=\max(f^*(uv)-f(uv),0)$, on a reversed arc $vu$,
$h(vu)=\max(f(uv)-f^*(uv),0)$. It respects the residual capacities, has the same net outflow, zero, at every vertex including $s$ and $t$ (the values are equal), and costs $\kappa(f^*)-\kappa(f)$. A nonnegative function with zero net outflow everywhere is a sum of positive multiples of
cycles of $D_f$: if $h\ne0$, walk along arcs with $h>0$; every vertex entered has an arc with $h>0$ leaving it, so the walk closes a cycle $C$; subtract $\min_Ch$ along $C$, which keeps the
net outflows zero and removes an arc from the support, and induct. Each cycle has nonnegative cost by hypothesis, so $\kappa(f^*)\ge\kappa(f)$.
\end{proof}

\begin{theorem}[successive shortest paths; Jewell, 1958; Busacker and Gowen, 1960]\label{thm:ssp}
Let $f$ be a minimum-cost flow of value $k$, and $P$ a cheapest path from $s$ to $t$ in $D_f$ (with respect to $\kappa$). Augmenting along $P$ by $\delta\le$ its bottleneck gives a minimum-cost flow of value $k+\delta$. Consequently, with nonnegative costs and integer capacities, starting from the zero flow
and repeating gives minimum-cost flows of every value up to the maximum; with Dijkstra's algorithm and Johnson's potentials, each step costs $O(m\log n)$.
\end{theorem}
\begin{proof}
By Lemma~\ref{lem:negcycle}, $D_f$ has no negative cycle, so the cheapest distances $\pi(v)$ from $s$ are well defined on the vertices reachable from $s$, and the \emph{reduced costs} $\kappa_\pi(uv)=\kappa(uv)+\pi(u)-\pi(v)$ are nonnegative on the arcs of $D_f$ between such vertices,
by the triangle inequality, and zero on the arcs of $P$. The new arcs of $D_{f'}$ are reverses of arcs of $P$, of reduced cost $0$ too. Let $C$ be a cycle in $D_{f'}$. If all its arcs are in $D_f$, its cost is nonnegative by hypothesis. Otherwise it uses a reversed arc of $P$, so it passes through
a vertex reachable from $s$ in $D_f$, and following its arcs (old arcs, or reversed arcs of $P$ whose ends are on $P$) it never leaves the reachable vertices; so its cost equals its reduced cost, which is nonnegative. By the lemma, $f'$ is optimal.

With nonnegative costs, the zero flow is optimal for value $0$, since $D$ itself has no negative cycle. With integer capacities each bottleneck is a positive integer. The reduced costs being nonnegative, the next cheapest path is found by Dijkstra's algorithm on $\kappa_\pi$
(Theorem~\ref{thm:johnson}), giving new distances and new potentials. The vertices reachable from $s$ in $D_{f'}$ were already reachable in $D_f$, since the only new arcs join vertices of $P$, so the potentials are
defined wherever Dijkstra's algorithm looks.
\end{proof}

The \emph{assignment problem} asks, for an $n\times n$ matrix of costs $c_{ij}$ (worker $i$ on job $j$), for a permutation $\sigma$ minimising $\sum_ic_{i\sigma(i)}$. It is the minimum-cost flow of value $n$ in the network $s\to$ workers $\to$ jobs $\to t$, all capacities $1$, costs $c_{ij}$ on the middle arcs
(an integral optimal flow exists, by the integrality of each augmentation).

\begin{corollary}[Kuhn, 1955; the Hungarian algorithm]\label{cor:hungarian}
The assignment problem is solved in time $O(n^3)$.
\end{corollary}
\begin{proof}
Theorem~\ref{thm:ssp} needs $n$ augmentations, each a Dijkstra computation on a graph with $O(n)$ vertices and $O(n^2)$ arcs; with an array instead of a heap, Dijkstra costs $O(n^2)$ on a dense graph (each of the $O(n)$ extractions scans the array, each arc is relaxed once).
Shifting all costs by a constant makes them nonnegative without changing the optimal permutations.
\end{proof}

Harold Kuhn named his 1955 method ``Hungarian'' because it rested on the work of D\'enes K\H onig and Jen\H o Egerv\'ary, whose 1931 paper he had translated from Hungarian. Half a century later, in 2006, Fran\c cois Ollivier wrote to him that an essentially identical algorithm was in
the posthumous papers of Carl Gustav Jacobi, who had died in 1851, written in Latin and motivated by systems of differential equations. Kuhn told the story in his 2010 lecture \emph{A tale of three eras}.

\begin{remark}
Flows remain a lively research area. Goldberg and Tarjan's \emph{push--relabel} method (1988), which does not maintain a flow but a \emph{preflow}, is the practical workhorse; Orlin (2013) reached $O(nm)$ for all networks; and in 2022 Chen, Kyng, Liu, Peng, Probst Gutenberg and Sachdeva
announced an algorithm computing maximum flows, and minimum-cost flows, in time $m^{1+o(1)}$, almost linear, by interior point methods with dynamic data structures, a striking meeting of continuous optimisation (Chapter~7) and combinatorics.
\end{remark}

\begin{beauty}
Every algorithm of this chapter improves a solution along a path and stops when no path is left, and every stopping point carries a proof of optimality: a cut whose capacity equals the flow, a set of odd vertices
that no matching can beat, potentials that make every cycle nonnegative. The algorithms are fast because the obstacles are understood: shortest paths make distances grow, phases make them grow faster, and Edmonds turned the one real obstacle, the odd
cycle, into a single vertex.
\end{beauty}

\section{Exercises}

\begin{exercise}[$\star$]\label{ex:ffbad}
In the network of Figure~\ref{fig:ffbad}, show that choosing augmenting paths badly takes $2C$ augmentations, while Edmonds--Karp takes two.
\end{exercise}

\begin{figure}[ht]
\centering
\fig{ch06-ffbad}
\caption{Four vertices, five arcs, maximum flow $2C$.}
\label{fig:ffbad}
\end{figure}

\begin{exercise}[$\star\star$]\label{ex:scaling}
\emph{Capacity scaling.} With integer capacities at most $U$, augment only along paths of bottleneck at least $\Delta$, for $\Delta=2^{\lfloor\log_2U\rfloor},\dots,2,1$, halving $\Delta$ when no such path remains. Show that each value of $\Delta$ allows at most $2m$ augmentations, hence time
$O(m^2\log U)$.
\end{exercise}

\begin{exercise}[$\star$]\label{ex:vertexcap}
Find a maximum number of internally vertex-disjoint paths from $s$ to $t$ in a directed graph with one flow computation, by splitting each vertex $v$ into $v^-\to v^+$ with capacity $1$.
\end{exercise}

\begin{exercise}[$\star\star$]\label{ex:dinicunit}
Show that on a network with all capacities $1$, a phase of Dinic's algorithm costs $O(m)$, and that there are $O(\sqrt m)$ phases. Deduce $O(m\sqrt m)$ in all.
\end{exercise}

\begin{exercise}[$\star$]\label{ex:konigcover}
Given a maximum matching of a bipartite graph, compute a vertex cover of the same size (K\H onig) in linear time.
\end{exercise}

\begin{exercise}[$\star$]\label{ex:whyblossom}
Run the alternating forest of Section~\ref{sec:blossom} on the opening example (triangle $xyz$, matching $\{yz\}$, pendant vertex $w$ at $z$). Where is the blossom, and what does the algorithm return?
\end{exercise}

\begin{exercise}[$\star$]\label{ex:tutteberge}
For a set $A$ of vertices, let $\mathrm{odd}(G-A)$ be the number of components of $G-A$ with an odd number of vertices. Show that every matching has at most $\frac12\big(n+|A|-\mathrm{odd}(G-A)\big)$ edges. (Equality for some $A$ is the Tutte--Berge formula.)
\end{exercise}

\begin{exercise}[$\star$]\label{ex:cyclecancel}
\emph{Cycle cancelling.} Show that with integer capacities and costs, repeatedly pushing flow around a negative cycle of the residual network, starting from a maximum flow, terminates with a minimum-cost maximum flow.
\end{exercise}

\begin{exercise}[$\star\star$]\label{ex:bottleneck}
\emph{Bottleneck assignment.} Find a perfect matching of a bipartite graph minimising the \emph{largest} cost of its edges, in time $O(m\sqrt n\log m)$.
\end{exercise}

\begin{exercise}[$\star\star$]\label{ex:edgeconn}
The \emph{edge connectivity} $\lambda(G)$ of a connected graph is the least number of edges whose removal disconnects it. Show that it can be computed with $n-1$ maximum flow computations.
\end{exercise}

\section{Hints}

\begin{hint}{ex:ffbad}
Alternate between $s\,a\,b\,t$ and $s\,b\,a\,t$, the second using the arc $ab$ backwards.
\end{hint}
\begin{hint}{ex:scaling}
When $\Delta$ is halved, the residual network has no path of bottleneck $\ge\Delta$: look at the set of vertices reachable by arcs of residual capacity $\ge\Delta$.
\end{hint}
\begin{hint}{ex:vertexcap}
Arcs $uv$ of the graph become $u^+v^-$ with capacity $1$ (or infinite).
\end{hint}
\begin{hint}{ex:dinicunit}
After $\sqrt m$ phases, $d(t)\ge\sqrt m$; decompose a maximum flow of the residual network into arc-disjoint paths.
\end{hint}
\begin{hint}{ex:konigcover}
Let $Z$ be the set of vertices reachable from free vertices of $X$ by alternating paths. Take $(X\setminus Z)\cup(Y\cap Z)$.
\end{hint}
\begin{hint}{ex:whyblossom}
The roots are $x$ and $w$. Which edge joins two even vertices first?
\end{hint}
\begin{hint}{ex:tutteberge}
Each odd component of $G-A$ contains a vertex that is free or matched into $A$.
\end{hint}
\begin{hint}{ex:cyclecancel}
The cost is an integer and is bounded below.
\end{hint}
\begin{hint}{ex:bottleneck}
Sort the costs and search for the least threshold $\theta$ such that the edges of cost $\le\theta$ contain a perfect matching.
\end{hint}
\begin{hint}{ex:edgeconn}
Fix a vertex $s$; some minimum cut separates $s$ from some $t$.
\end{hint}

\section{Solutions}

\begin{solution}{ex:ffbad}
\idea A bad path uses the middle arc, alternately forwards and backwards, and gains $1$ each time.
\writeup Augmenting along $s\,a\,b\,t$ (bottleneck $1$, the arc $ab$) and then along $s\,b\,a\,t$ (bottleneck $1$, the arc $ab$ reversed) increases the value by $2$ and returns the middle arc to $0$; repeating $C$ times reaches the maximum $2C$ after $2C$ augmentations,
which is exponential in the size of the input (about $\log C$ bits). Edmonds--Karp chooses the shortest paths $s\,a\,t$ and $s\,b\,t$, of length $2$, each with bottleneck $C$, and is done after two augmentations.
\end{solution}

\begin{solution}{ex:scaling}
\idea At the start of a scaling phase the remaining gap is small, and each augmentation of the phase closes a large part of it.
\writeup At the end of the phase with parameter $2\Delta$ (or at the start, for the first one, where the maximum flow is at most $mU<2m\Delta$), let $S$ be the set of vertices reachable from $s$ by residual arcs of capacity $\ge2\Delta$; $t\notin S$. Every residual arc leaving $S$ has capacity $<2\Delta$, and
there are at most $m$ of them (one per arc of $D$, as arcs are not opposite), so by the argument of Tome~C1, Theorem~6.6(1), the maximum flow exceeds the current one by less than $2m\Delta$. Each augmentation of the phase $\Delta$ adds at least $\Delta$, so there are fewer than $2m$ of them.
There are $\lfloor\log_2U\rfloor+1$ phases, and each augmentation costs $O(m)$ (a search in the arcs of residual capacity $\ge\Delta$). At the end, with $\Delta=1$, no augmenting path remains, and the flow is maximum.
\end{solution}

\begin{solution}{ex:vertexcap}
\idea Turning each vertex into an arc of capacity $1$ makes ``vertex-disjoint'' into ``arc-disjoint''.
\writeup Build $D'$ with vertices $v^-,v^+$, arcs $v^-v^+$ of capacity $1$ for $v\ne s,t$, and arcs $u^+v^-$ of capacity $1$ for each arc $uv$ (with $s^+=s$, $t^-=t$). An integral flow of value $k$ decomposes into $k$ paths (Tome~C1, Lemma~6.7) using each $v^-v^+$ at most once, which become $k$ internally vertex-disjoint
paths of $D$; conversely such paths give a flow. So the maximum flow is the answer (and its minimum cut is Menger's separating set).
\end{solution}

\begin{solution}{ex:dinicunit}
\idea With unit capacities, every arc on an augmenting path is saturated, and long distances leave little room for flow.
\writeup In a phase, every augmentation of the blocking flow saturates every arc of its path, which is then deleted; so each arc is advanced over at most once before being deleted or retreated over, and the phase costs $O(m)$. After $\sqrt m$ phases,
$d_f(t)\ge\sqrt m$. The residual network $D_f$ is again a unit network (forward and reversed arcs have capacity $1$), and a maximum flow in it, which exists with value $|f^*|-|f|$, decomposes into arc-disjoint paths from $s$ to $t$, each with at least $\sqrt m$ arcs. Since $D_f$ has at most $m$ arcs that can be used
(one per arc of $D$), there are at most $\sqrt m$ such paths: at most $\sqrt m$ units of flow are missing, and each further phase adds at least one. So there are at most $2\sqrt m$ phases.
\end{solution}

\begin{solution}{ex:konigcover}
\idea The search that fails to find an augmenting path draws the cover: what it reached on the $Y$ side, and what it missed on the $X$ side.
\writeup Let $M$ be maximum and $Z$ the set of vertices reachable from free vertices of $X$ by $M$-alternating paths (computed by one search). Let $K=(X\setminus Z)\cup(Y\cap Z)$. Alternating paths from free vertices of $X$ go from $X$ to $Y$ by non-matching edges and from $Y$ to $X$ by matching edges. Let $xy$ be an edge with $x\in X\cap Z$. If $xy\in M$, then $x$ is matched, so it is not a starting
vertex, and an alternating path reaches it through its mate $y$: $y\in Z$. If $xy\notin M$, an alternating path to $x$ extended by $xy$ reaches $y$: $y\in Z$. So every edge meets $K$. Every $y\in Y\cap Z$ is matched (otherwise an augmenting path), with its mate in $Z$; every $x\in X\setminus Z$ is matched (free vertices of $X$ are in $Z$), with its mate
not in $Z$ (else $x$ would be). So each edge of $M$ contains exactly one vertex of $K$, and $|K|=|M|$. This is the proof of Tome~C1, Theorem~6.4, read as an algorithm.
\end{solution}

\begin{solution}{ex:whyblossom}
\idea The triangle is a blossom whose base is the root itself.
\writeup The roots are the free vertices $x$ and $w$. From $x$ (even), the edge $xy$ grows $y$ (odd) and its mate $z$ (even). Scanning from $x$ the edge $xz$, or from $z$ the edge $zx$, joins two even vertices of the same tree: the blossom $B=x\,y\,z$ with base $x$, and an empty stem.
Shrinking gives the edge $Bw$ between two free vertices: the recursive call returns the matching $\{Bw\}$, and expanding it matches $z$ (the vertex of $B$ joined to $w$) with $w$ and the remaining path $x\,y$ by the edge $xy$: the perfect matching $\{zw,xy\}$. (Had $z$ scanned
the edge $zw$ first, the forest would have found case~2 directly: $w$ is even, in another tree.)
\end{solution}

\begin{solution}{ex:tutteberge}
\idea Odd components cannot be perfectly matched inside themselves; each needs help from $A$, and $A$ can help at most $|A|$ of them.
\writeup Let $N$ be a matching. In each odd component $K$ of $G-A$, the edges of $N$ inside $K$ cover an even number of vertices, so some vertex of $K$ is free in $N$ or matched to a vertex outside $K$, which must be in $A$. At most $|A|$ components are helped by $A$, so at least
$\mathrm{odd}(G-A)-|A|$ vertices are free, and $|N|\le\frac12\big(n-\mathrm{odd}(G-A)+|A|\big)$.
\end{solution}

\begin{solution}{ex:cyclecancel}
\idea Each cancellation lowers an integer cost, which cannot fall forever.
\writeup Pushing the bottleneck around a negative cycle keeps the flow feasible and of the same value, and with integer data the bottleneck is a positive integer and the cost decreases by at least $1$. The cost of any flow is at least $-\sum_a|\kappa(a)|c(a)$, so the process stops,
and when it stops there is no negative cycle: by Lemma~\ref{lem:negcycle} the flow has minimum cost among maximum flows. (With the cycle of least \emph{mean} cost each time, Goldberg and Tarjan made this strongly polynomial.)
\end{solution}

\begin{solution}{ex:bottleneck}
\idea ``Is there a perfect matching using only edges of cost at most $\theta$?'' is monotone in $\theta$: binary search.
\writeup Sort the distinct costs $\theta_1<\dots<\theta_r$, $r\le m$. The predicate ``the edges of cost $\le\theta_i$ contain a perfect matching'' is false then true as $i$ grows; the answer is the first $\theta_i$ where it holds. Binary search tests it $O(\log m)$ times, each by Hopcroft--Karp
in $O(m\sqrt n)$; sorting costs $O(m\log m)$.
\end{solution}

\begin{solution}{ex:edgeconn}
\idea A minimum cut separates our fixed vertex from someone, and the max-flow min-cut theorem measures each separation.
\writeup Replace each edge by two opposite arcs of capacity $1$. For vertices $s\ne t$, the maximum flow $\lambda(s,t)$ is the least number of edges separating $t$ from $s$ (max-flow min-cut, Tome~C1, Theorem~6.6, and Menger). Fix $s$. A minimum disconnecting set of edges splits $V$ into two sides,
and some vertex $t$ lies on the side not containing $s$; so $\lambda(G)=\min_{t\ne s}\lambda(s,t)$: $n-1$ flow computations. (Stoer and Wagner gave in 1997 a simpler algorithm without flows.)
\end{solution}

\section*{Sources and further reading}

\begin{itemize}
\item \emph{The stories.} J.~Edmonds, \emph{Paths, trees, and flowers}, Canadian Journal of Mathematics 17 (1965), 449--467, Section~2 for the ``good algorithm'' digression; H.~W.~Kuhn, \emph{The Hungarian method for the assignment problem}, Naval Research Logistics Quarterly 2 (1955), 83--97;
H.~W.~Kuhn, \emph{A tale of three eras: the discovery and rediscovery of the Hungarian Method}, European Journal of Operational Research 219 (2012), 641--651 (a plenary lecture of 2010), for Jacobi and Ollivier; Y.~Dinitz, \emph{Dinitz' algorithm: the original version and Even's version}, in
\emph{Theoretical Computer Science: Essays in Memory of Shimon Even}, Lecture Notes in Computer Science 3895, Springer, 2006, 218--240.
\item J.~Edmonds and R.~M.~Karp, \emph{Theoretical improvements in algorithmic efficiency for network flow problems}, Journal of the ACM 19 (1972), 248--264; E.~A.~Dinic, \emph{Algorithm for solution of a problem of maximum flow in a network with power estimation}, Soviet Mathematics Doklady 11 (1970),
1277--1280; J.~E.~Hopcroft and R.~M.~Karp, \emph{An $n^{5/2}$ algorithm for maximum matchings in bipartite graphs}, SIAM Journal on Computing 2 (1973), 225--231; U.~Zwick, \emph{The smallest networks on which the Ford--Fulkerson maximum flow procedure may fail to terminate}, Theoretical Computer Science 148 (1995), 165--170.
\item S.~Micali and V.~V.~Vazirani, \emph{An $O(\sqrt{|V|}\,|E|)$ algorithm for finding maximum matching in general graphs}, Proceedings of the 21st IEEE FOCS (1980), 17--27; A.~Schrijver, \emph{Combinatorial Optimization: Polyhedra and Efficiency}, Springer, 2003, Chapters~9--12 and 24, for the Tutte--Berge
certificate from Edmonds' algorithm and much more.
\item W.~S.~Jewell, \emph{Optimal flow through networks}, Interim Technical Report 8, MIT Operations Research Center, 1958; R.~G.~Busacker and P.~J.~Gowen, \emph{A procedure for determining a family of minimum-cost network flow patterns}, ORO Technical Paper 15, 1960; N.~Tomizawa, \emph{On some techniques useful for solution
of transportation network problems}, Networks 1 (1971), 173--194.
\item A.~V.~Goldberg and R.~E.~Tarjan, \emph{A new approach to the maximum-flow problem}, Journal of the ACM 35 (1988), 921--940; J.~B.~Orlin, \emph{Max flows in $O(nm)$ time, or better}, Proceedings of the 45th ACM STOC (2013), 765--774; L.~Chen, R.~Kyng, Y.~P.~Liu, R.~Peng, M.~Probst Gutenberg and S.~Sachdeva,
\emph{Maximum flow and minimum-cost flow in almost-linear time}, Proceedings of the 63rd IEEE FOCS (2022), 612--623; M.~Stoer and F.~Wagner, \emph{A simple min-cut algorithm}, Journal of the ACM 44 (1997), 585--591.
\end{itemize}
